\section{Market state, latent fields, and valid events}\label{sec:state}
\subsection{Displayed state and the role of marks}
Fix a finite physical price grid and a finite collection of admissible order sizes. A displayed state $x$ records nonnegative integer queue sizes, relevant order identifiers if available, and any public information needed by the transition rule. A full latent state $z$ can contain reserve quantities and other predictive variables. The observed history $\Hist_{t-}$ includes the initial state, prior event times and marks, and the policy's previous actions. The filtration generated by this history is denoted $\F^O_{t-}$.

Let $\mathcal E$ be a finite mark set. A mark contains enough information that the displayed update is a deterministic map
\begin{equation}\label{eq:transition}
 x_t=\Delta(x_{t-},e_t).
\end{equation}
For example, a combined execution-and-replenishment mark must specify both executed displayed quantity and the resulting displayed replenishment. A mark that merely records an execution need not determine the next displayed state when reserves are hidden. Splitting an execution and its replenishment into two linked marks is an alternative representation, provided their timing convention is explicit.

Let $\mathcal V$ be the set of states satisfying the simulator's stated market rules. For each $x\in\mathcal V$, let $\mathcal E(x)$ contain only marks for which $\Delta(x,e)\in\mathcal V$. A simple continuous-trading model can require integer lots, nonnegative queues, cancellation quantities no larger than the relevant order, and a correctly updated best bid and ask. A limit order that crosses the opposite quote is handled by the specified matching rule. These are model requirements; a production implementation needs its own venue and feed specification. Auctions require a separate transition regime and are not represented by an unmodified continuous-trading rule.

\begin{proposition}[Structural event validity]\label{prop:validity}
Suppose $x_0\in\mathcal V$, the event process is nonexplosive on $[0,T]$, its predictable intensities are nonnegative, and the intensity of every $e\notin\mathcal E(x_{t-})$ is zero. Then $x_t\in\mathcal V$ for all $t\le T$ almost surely.
\end{proposition}
\begin{proof}
On a nonexplosive path there are finitely many jumps before $T$. Between jumps the displayed state is constant. At each jump, the sampled mark belongs to $\mathcal E(x_{t-})$ almost surely, so its transition maps a valid state to another valid state. Induction over the jump times proves the assertion.
\end{proof}

A penalty in a training loss does not imply this proposition. The support restriction belongs to the sampler and transition mechanism. Conversely, structural validity does not imply that the rates are statistically accurate.

\subsection{Reserve accounting}
At a fixed price, write $D,H\in\N$ for displayed and hidden lots. An execution consumes $q\le D$ displayed lots, and a replenishment transfers $r\le H$ hidden lots into display:
\begin{equation}\label{eq:reserve-update}
 D'=D-q+r,\qquad H'=H-r.
\end{equation}
Additional restrictions on $r$, such as replenishment only after an order's displayed peak is exhausted, are imposed by the specified admissible-mark set. They do not alter the following accounting identity.

\begin{proposition}[Execution-linked reserve conservation]\label{prop:reserve}
Under \eqref{eq:reserve-update},
\begin{equation}\label{eq:reserve-conservation}
 D'+H'=D+H-q,\qquad D',H'\ge0.
\end{equation}
For a sequence of such events with no new submissions or cancellations,
$D_n+H_n=D_0+H_0-\sum_{k<n}q_k$.
\end{proposition}
\begin{proof}
The transfer $r$ cancels when the two update equations are added. Nonnegativity follows from $q\le D$ and $r\le H$. Summing the one-step identity telescopes.
\end{proof}

Thus replenishment is a transfer, not new aggregate liquidity. Hidden submissions, cancellations, and venue-specific refresh rules can be added as separate transitions. A continuous latent field may parameterize their rates, but it does not replace integer inventory accounting. If a latent decoder rounds a real coordinate to an integer reserve, that rounding is discontinuous; one must verify the sensitivity of the resulting intensity map rather than inherit it from the pre-decoder field.

\subsection{Field-valued rates and numerical resolution}
Let $(A,\nu)$ be a finite measure space, with $M=\nu(A)<\infty$. Its coordinates may encode price offset, event type, and size. A nonnegative field
\[
 f_t(\Hist,z,\cdot)\in L^2(A,\nu)
\]
is an intensity \emph{density} with respect to $\nu$. Fix disjoint cells $(B_e)_{e\in\mathcal E}$ representing physical marks. With an observed-state feasibility mask $m_e(\Hist)\in\{0,1\}$, define
\begin{equation}\label{eq:cell-rates}
 \lambda_e(t,\Hist,z,a)
 =m_e(\Hist,a)\int_{B_e}f_t(\Hist,z,a,u)\,\nu(\dd u).
\end{equation}
The cells need not cover all of $A$; unused regions carry no sampled marks. If an admissibility condition depends on $z$, include it in the conditional intensity field and verify the required latent sensitivity for that complete field. A common observed-state mask alone does not encode every hidden-inventory constraint.

The representation in \eqref{eq:cell-rates} matters. Evaluating an $L^2$ function at an arbitrary point is not a continuous operation, and an $L^2$ equivalence class has no canonical point values. Integrating over physical cells is well defined and gives the partition-independent estimate used below. Refining numerical quadrature inside those cells leaves the physical market specification fixed. Changing the actual tick size, allowed order sizes, or matching rule changes the modeled market.
